\documentclass{article}
\usepackage{hyperref}

\begin{document}
\subsection*{7.2 Unidades}

\noindent\textbf{Amperio [A]:} unidad de corriente el\'ectrica, representando el caudal de carga el\'ectrica.
\newline
\noindent\textbf{Carga electr\'onica[e]:} magnitud de carga el\'ectrica transportada por un electr\'on.
\newline
\noindent\textbf{Fotogramas por segundo [fps]:} cantidad de im\'agenes consecutivas tomadas por la c\'amara en un segundo.

\noindent\textbf{Candela [cd]:} Intensidad luminosa.

\noindent\textbf{Decibelio [dB]:} Unidad logar\'itmica usada para medir la relaci\'on entre dos magnitudes de potencia.

\noindent\textbf{Lumen [lm]:} Unidad SI de flujo luminoso que describe la luz producida por una fuente por unidad de tiempo. Esta unidad es com\'un en iluminaci\'on continua y bombillas de flash.

\noindent\textbf{Lumen-segundo[lm-s]:} Unidad SI de luz producida a lo largo del tiempo, o energ\'i luminosa. Esta unidad es menos com\'un, pero usada por bombillas de flash.

\noindent\textbf{Lux [lx]:} Unidad SI de illuminance? o flujo luminoso por unidad de \'area, igual a un lumen por metro cuadrado. Lux es comunmente usado como una medida de intensidad de luz (o flujo luminoso). Esta unidad describe la cantidad de luz incidente en un objeto por unidad de \'area y es mas comun para iluminaci\'on continua, pero puede ser utilizado para describir luz estrobosc\'opica tambi\'en.

\noindent\textbf{Pixel [px]:} V\'ease definiciones en el glosario para *pixel de imagen* \'o *pixel de detector*. Se usan prefijos SI, i.e. Mpx, Gpx, etc.

\noindent\textbf{Pixels por segundo [px/s]:} Ratio de transferencia de datos desde el detector digital de la c\'amara hacia la memoria de la c\'amara. Se usan prefijos SI, e.g. Gpx/s.

\noindent\textbf{Puntos o muestras por segundo, Hertzio [Hz]:} frequencia de muestreo.

\noindent\textbf{Voltio [V]:} Unidad standard de potencial el\'ectrico.

\noindent\textbf{Vatio [W]:} Unidad standard de potencia o flujo radiante.

\noindent\textbf{Vatio-segundo [W-s] \'o Julio [J]:} Unidad comunmente usada para la potencia producida sobre un periodo de tiempo. Esta unidad es comun para iluminaci\'on estrobosc\'opica/pulsada.


\end{document}
